\documentclass[final,a4paper,10pt]{report}
\def\RuleAssetPath{rules/2026/formula-competition/assets}
\input{template}

\begin{document}
\thispagestyle{firstpage}

\begin{center}
  {\pretendardb\fontsize{24}{28}\selectfont\addfontfeatures{LetterSpace=5} KSAE 대학생 자작자동차대회}\\[2ex]
  {\pretendardb\fontsize{24}{28}\selectfont\addfontfeatures{LetterSpace=5} Formula 경기진행규정}\\[2.5ex]
  [ 시행 2026.3.19. 이사회 ]
\end{center}

\chapter{목적 및 일반사항}

\section{목적}
본 규정은 대학생 자작자동차대회 대회운영규정(이하 ``대회운영규정''이라 한다) 제10조 4항에 따라 Formula 경기진행에 관한 사항을 규정함을 목적으로 한다.

\section{일반사항}
Formula 경기진행의 일반적인 사항들은 대회운영규정에 따르며, Formula 제작에 관한 사항은 Formula 차량기술규정에 따른다.

\chapter{차량검사 진행}

\section{차량검사}\label{rule:formula-competition.vehicle-inspection}
\begin{enumerate}
  \item 차량검사의 목적은 차량이 규정을 만족하는 설계와 안전 요구사항, 취지에 맞게 제작되었는지를 확인하기 위함이다.
  \item 지정된 시각까지 차량검사를 통과하지 못한 차량은 조직위원회에 의해 특별조치가 인정되지 않는 한 경기에 출전할 수 없다.
  \item 조직위원회는 안전에 문제가 있는 차량의 출전을 금지할 수 있다.
  \item 차량검사 대상 팀은 차량검사의 모든 준비를 완료한 후 지정된 검차장소에 차량을 준비한다.
  \item 차량검사 진행
    \begin{enumerate}
      \item 경기에 참가하기 위한 허가를 받거나 트랙에서 주행하기 전에, 각 차량은 차량검사의 모든 부분을 통과해야 한다. 정확한 검사와 시험을 위해 장비가 사용될 수 있으며, 차량검사의 진행에 관한 권한은 조직위원회가 가진다.
      \item 조직위원회는 조사 양식에 포함된 모든 항목에 따라 차량의 안전과 기능에 관한 조사를 진행하며, 조직위원회가 규정 적합여부 검사를 추가로 원하는 경우 다른 항목을 포함하여 조사를 수행할 수 있다.
      \item 차량검사를 통과한 차량은 경기기간 동안 ``검차 통과'' 상태를 유지하여야 하며, 임의로 수정해서는 안 된다.
    \end{enumerate}
  \item 차량 수정과 재검\\
    참가팀은 차량의 특정 부분이 규정에 벗어나거나 안전하지 않다고 지적된 사항에 대해서 해당 사항을 수정하고 차량 재검사를 받아야 한다.
    \begin{enumerate}
      \item 조직위원회는 대회기간 중 언제든지 임의의 차량을 재조사할 수 있고, 부적합한 사항에 대하여 수정을 요구할 수 있다.
      \item 검차 운영시간을 초과하여 검차를 통과하지 못한 차량은 실격 처리되고 모든 경기 참가가 제한된다.
    \end{enumerate}
  \item 차량검사의 완료\\
    안전검사(전기동력차량 전기기술검사 포함), 틸팅검사, 제동검사, 소음검사(기관동력) / 음향검사(전기동력) / 우천검사(전기동력) 등 각 검사별로 통과 스티커가 제공되며, 전부 통과하여 스티커가 모두 부착된 차량만이 동력계 시동, 연습 주행 및 경기의 참가가 허가된다.
\end{enumerate}

% This tail is declared here and emitted after Articles 4--10 below. Keeping the
% source order explicit through one macro avoids duplicating the long official text.
\newcommand{\CompetitionRulesTail}{%
\section{내구레이싱 경기}
\begin{enumerate}
  \item 내구레이싱 경기는 차량의 전체 성능과 신뢰성을 테스트하여 차량의 성능에 대한 올바른 판단을 하고자 함에 있다.
  \item 내구레이싱 경기코스
    \begin{enumerate}
      \item 내구레이싱 경기코스는 직선 구간, 회전 구간, 헤어핀 구간, 슬라럼 구간, 시케인 구간 등을 포함한다.
      \item 내구레이싱 경기코스의 표준사항은 다음과 같다.
        \begin{itemize}
          \item 직선 구간: 양쪽 끝에 헤어핀을 가진 최대 70m 구간 또는 끝부분에 완만한 곡선을 가진 60m를 넘지 않는 구간. 직선구간에는 추월구간이 포함되어 있을 수 있다.
          \item 회전 구간: 30m에서 60m의 직경을 가진 구간
          \item 헤어핀 구간: 최소 9m의 외부 직경을 가진 구간
          \item 슬라럼(지그재그) 구간: 7m에서 15m 사이의 직선으로 놓여진 파일런 구간
          \item 기타 구간: 시케인(자동차 경주도로에서 감속을 위한 장애물), 다양한 회전, 반경이 줄어드는 회전 구간 등을 가지며 트랙의 최소 폭은 4m로 만들어진다.
        \end{itemize}
      \item 경기장 상황에 따라 내구레이싱 경기코스에 단차, 배수홀 등의 구간이 포함될 수 있으며, 해당 구간은 경기장 특성으로 참가팀들은 이에 대비한 차량 세팅과 주행에 임해야 한다. 노면 상황에 의한 차량 파손은 참가팀에 책임이 있다.
    \end{enumerate}
    \fig{내구레이싱 경기 코스 예}{competition}{0.72}
  \item 내구레이싱 경기의 진행
    \begin{enumerate}
      \item 내구레이싱 경기의 총 주행거리는 조직위원회에서 약 20km 전후로 결정하며 경기 진행 중 참가팀의 허용되지 않은 임의 차량 조작은 허용되지 않는다.
      \item 내구레이싱 경기 중 드라이버 교체가 요구된다.
      \item 차량의 추월은 추월구간 또는 코스 경기진행요원의 지시 하에서만 가능하다.
      \item 조직위원회 결정에 따라 그룹 A와 그룹 B 경기를 구분하여 진행할 수 있다.
      \item 내구레이스 경기 차량은 주어진 위치에 주행시간 계측기(트랜스폰더) 설치 의무가 있으며 경기 상황에 따라 계측기 대신 스톱워치를 이용하여 주행시간을 계측할 수 있다.
    \end{enumerate}
  \item 내구레이싱 경기의 출발 순서
    \begin{enumerate}
      \item 내구레이싱 경기의 출발 순서는 오토크로스 경기 기록으로 정해지며, 만약 오토크로스 경기의 기록이 없는 차량의 경우 가속경기, 스키드패드 경기의 기록 순으로 정하도록 한다. 만약 어떠한 기록도 없는 차량의 경우 조직위원회는 마지막 순서로 내구레이싱 경기에 임할 수 있도록 차량번호 순으로 정해 실시한다.
      \item 조직위원회에서는 종목별 평가 결과와 차량의 주행조건 등을 고려하여 차량의 출발순서를 조정할 수 있다.
      \item 그룹 A와 그룹 B의 경기진행 순서는 조직위원회에서 결정한다.
      \item 참가팀은 내구레이싱 경기를 위해 급유 또는 배터리 충전 등 준비를 완료하고 차량의 출발 순서에 따라 대기하여야 한다.
      \item 자신의 트랙 진입 순서에 출발할 준비가 되어 있지 않은 차량은 경기 시간이 허락되면 조직위원회의 결정에 따라 마지막 순서에 진행될 수 있으며 진입 순서 후행 횟수마다 2분의 벌칙 시간이 주어진다.
    \end{enumerate}
  \item 차량의 급유 또는 차량의 충전
    \begin{enumerate}
      \item 내구레이싱 경기에 참가하기 전에 각 차량은 급유 및 배터리 충전을 마쳐야 하며, 내구경기 진행 중 급유 및 충전은 할 수 없다.
      \item 내구경기 참가 전 차량의 급유량, 충전량 등을 확인을 요청할 수 있다.
    \end{enumerate}
  \item 내구레이싱 경기의 출발 및 재출발
    \begin{enumerate}
      \item 경기가 시작되면 차량은 언제라도 외부의 도움 없이 출발 및 재출발이 가능하다.
      \item 차량이 트랙에서 정지하여 벗어난다면 재출발까지 1분의 시간이 허락된다. 또한 코스에서 주행이 불가능할 경우에는 운영요원에 의해 코스 밖으로 차량이 이동될 수 있으며, 재시동까지 최대 2분의 시간이 주어지고 2분 초과 시에는 DNF 처리한다.
      \item 드라이버 교체 후 차량의 출발에 문제가 발생한다면 재출발까지 최대 2분의 시간이 주어지며, 주어진 시간에 재출발이 불가능한 차량은 운행불가 차량으로 판단하여 DNF 처리한다.
    \end{enumerate}
  \item 내구레이싱 경기의 드라이버 교체 절차
    \begin{enumerate}
      \item 첫 번째 드라이버는 주어진 주행거리의 1/2을 주행하며 드라이버 교체 사인을 받으면 드라이버 교체 구역으로 들어와야 한다.
      \item 두 번째 드라이버는 잔여 주행거리를 주행하며, 내구레이스 경기를 마치게 될 때 주행시간이 기록된다.
      \item 드라이버 교체를 위해 3분의 시간이 주어지며, 3분이 경과된 후부터 초과된 시간은 전체 주행기록시간에 추가된다.
      \item 드라이버 교체를 위한 지정된 공간에는 팀당 교체할 드라이버를 포함하여 3명의 인원만이 출입 가능하다. 전기동력 차량 팀의 경우 3명 중 한 명은 반드시 ESO이어야 한다.
      \item 두 번째 드라이버를 위한 차량조정이나 날씨에 따른 타이어 교체를 위한 공구만 휴대할 수 있고, 조정작업은 드라이버 승하차 전후에만 가능하며 소모품 교환은 허락되지 않는다.
      \item 드라이버 교체 구역에 다른 팀원이 출입할 경우 내구레이싱 경기 최종점수에서 1인당 50점이 감점된다.
    \end{enumerate}
  \item 트랙 입장
    \begin{enumerate}
      \item 코스 내 주행차량의 밀집정도에 따라 조직위원회의 결정으로 차량의 트랙입장이 허가된다.
      \item 코스 내 주행 가능한 차량의 수는 트랙의 길이, 형태 그리고 경기 진행 상태에 따라 조직위원회에서 결정한다.
      \item 트랙에 입장한 차량의 숫자는 드라이버 교체지역에 있는 차량도 포함한다.
    \end{enumerate}
  \item 차량고장과 운행정지
    \begin{enumerate}
      \item 차량고장 발생 시 차량은 코스에서 이동조치 되며, 재출발은 허락되지 않는다.
      \item 차량 정지 후 자가 출발이 가능하거나 차량과 지면사이에 파일런이 끼는 등의 경우에는 재출발이 가능하며, 코스를 벗어난 곳에서 재진입하여야 하며, 차량의 정비는 허락되지 않는다.
      \item 차량이 정지하여 정비나 외부의 도움 없이 차량이 재출발이 불가능할 경우, 경기진행요원의 지시에 따라 두 명의 팀원이 차량을 트랙 밖으로 이동시켜야 한다.
    \end{enumerate}
  \item 내구레이싱 경기의 최소 속도규정
    \begin{enumerate}
      \item 차량의 주행속도가 느려 코스에 주행 중인 다른 차량의 주행에 방해가 된다고 판단될 경우 즉시 코스 아웃 및 DNF 지시를 받을 수 있다.
      \item 최소 속도규정의 적용은 조직위원회의 결정에 따른다.
    \end{enumerate}
  \item 내구레이싱 경기 랩타임\\
    내구레이싱 경기 랩타임은 각 차량에 장착된 전자장비 또는 경기진행요원에 의해 직접 측정되며, 주행기록시간은 총 시간에서 별도의 드라이버 교체시간을 빼고, 벌칙으로 부가된 시간을 더하여 결정된다.
  \item 내구레이싱 경기 벌칙
    \begin{enumerate}
      \item 사고회피나 충분히 납득할 만한 이유라고 인정되면 벌칙이 부과되지 않는다.
      \item 파일런이 넘어지거나 움직였을 경우\\
        파일런 하나당 2초의 벌칙시간이 더해진다. (출발선과 결승선 후의 파일런도 포함)
      \item 코스 이탈\\
        코스 이탈은 차량의 네 바퀴 모두가 파일런이나 코스 선 등 정해진 코스의 경계선을 넘어설 경우를 말한다. 코스를 이탈하게 되면 드라이버는 이탈 발생지점이나 발생지점 전에서 코스로 재진입하여야 하며, 이를 위반 시 20초의 벌칙시간이 더해진다. 코스 이탈 중 지나친 파일런 하나당 2초의 벌칙 시간이 추가로 부과된다. 코스에서 두 바퀴가 이탈할 경우 즉각적인 벌칙은 부여되지 않으나, 반복적으로 발생할 경우 흑색기가 발효될 수 있다. 사고를 피하기 위해 또는 경기진행요원의 요구에 의한 코스 이탈은 벌칙이 부과되지 않는다.
      \item 슬라럼을 놓친 경우\\
        주어진 슬라럼 파일런의 하나 또는 그 이상의 게이트를 놓친 경우 각각 하나의 코스 이탈로 계산되어 각 20초의 벌칙시간이 부과된다.
      \item 주행 중 위반에 대한 벌칙\\
        다음과 같은 경우 벌칙이 부과되며 실격될 수 있다.
        \begin{enumerate}
          \item 깃발 불복종: 1분 추가
          \item Over Driving(흑색기에 가까운 경우 시): 1분 추가
          \item 다른 차량에 대한 접촉: 고의성 여부에 따라 2분 추가 또는 실격(조직위원회 판단)
        \end{enumerate}
      \item 규정위반 주행 시(주어진 코스를 고의적으로 이탈하여 주행할 경우): 2분 추가
      \item 차량의 기계적 문제\\
        경기 중 차량 문제점 및 안전 문제 발생이 의심될 경우 차량을 드라이버 교체 장소나 추월 위치로 진입시켜 이상 유무 확인을 진행할 수 있으며, 주행이 불가하다고 판단될 경우 DNF 처리될 수 있다.
      \item 공격적인 주행\\
        무모하고 공격적인 주행행동(물리적으로 다른 차를 트랙에서 벗어나게 하거나, 차량의 추월을 방해, 차 접촉의 가능성을 일으키게 하는 근접주행), 의도적인 코스 이탈을 하는 드라이버에게 흑색기를 발효한다. 드라이버가 흑색기 신호를 받으면 드라이버는 벌칙 구간으로 이동해야 하며 벌칙 내용에 대하여 확인할 수 있다. 벌칙 구간에서 벌칙내용을 확인하기 위해 소요된 시간은 주행시간에 포함된다.
      \item 깃발이 발령되고 발령된 랩을 포함하여 2랩을 초과하여 주행할 경우 실격(DNF) 처리한다.
    \end{enumerate}
  \item 내구레이싱 경기 점수계산
    \begin{enumerate}
      \item 내구레이싱 경기의 점수계산은 다음과 같다.
        \[
          \text{내구레이싱 경기 점수} = 250 \times
          \frac{((1.45 \times T_{\min}) / T_{\mathrm{yours}}) - 1}
               {((1.45 \times T_{\min}) / T_{\min}) - 1} + 100
        \]
        $T_{\mathrm{yours}}$는 해당 팀의 주행시간 합계, $T_{\min}$은 가장 빠른 팀의 최소 주행시간이다.
      \item 가장 빠른 내구경기 기록의 145\%를 초과한 차량은 내구레이싱 점수 중 주행성능에 대한 점수는 부여하지 않고 100점의 완주점수만 주어진다. 완주하지 못한 차량에 대해서는 DNF 처리한다.
    \end{enumerate}
  \item 동력계 확인\\
    조직위원회에서 정한 방법에 따라 경기 후 내연기관 차량의 엔진상태와 흡기제한장치 등 동력계통과 전기차량의 모터와 배터리 등 전기 동력계통을 즉시 확인할 수 있다.
  \item 내구레이싱 경기에서의 주행
    \begin{enumerate}
      \item 드라이버는 내구레이싱 경기 동안 여러 대의 차량이 코스에서 동시에 주행하므로 모든 규정과 드라이버 요구사항을 엄격히 따라야 한다.
      \item 공격적인 주행, 지시 불복종, 추월구간 양보 불복종 등의 경우 흑색기가 발령된다. 드라이버는 정해진 벌칙 구간으로 차량을 대기해야 하며 코스 경기진행요원에 의해 벌칙 시간이 부가된다. 벌칙 구간에서 대기시간은 경기진행요원에 의해 0분에서 4분까지 정해지며 주행시간에 포함된다.
      \item 주행 중 차량의 이상이 의심될 경우 흑색기가 발령되며, 정해진 장소에서 차량의 상태를 점검받아야 하며 대기 중 시간은 주행시간에 포함된다.
    \end{enumerate}
  \item 내구레이싱 경기에서의 추월
    \begin{enumerate}
      \item 내구레이싱 경기가 진행되는 동안 추월은 경기진행요원의 지시 하에 지정된 추월구간에서만 가능하다.
      \item 추월구간은 추월차선과 우회차선의 평행한 2차선으로 구성되며, 추월 차선을 이용하여 우회 차선의 차량을 추월한다. 추월구간에 접근하였을 때 주행이 느린 차량은 청색기를 받게 되며, 우회 차선으로 입구부터 감속하며 진입 후 지시 전까지 정차하며 대기해야 한다. 뒤따르는 빠른 차량은 기존차선을 계속 주행하여 추월한다. 우회 차선으로 진입한 차량은 추월구간 끝에서 주어지는 깃발의 지시를 따라 주행구간으로 재진입해야 한다.
      \item 우회 차선의 위치는 코스의 상황에 따라 우측이나 좌측에 위치할 수 있다.
      \item 운행이 불가능한 차량이나 트랙을 벗어난 차량을 추월할 경우 드라이버는 주행속도를 줄이고 모든 차량과 경기진행요원의 지시를 따르며 운전하여야 한다.
      \item 일반적인 주행에서는 모든 차량들은 추월차선을 이용하여 주행한다.
    \end{enumerate}
  \item 코스워크\\
    내구레이싱 경기 코스는 내구레이싱 경기 전에 조직위원회의 허가 하에 드라이버들이 도보로 살펴볼 수 있으며, 모든 드라이버는 내구레이싱 경기 전에 코스를 반드시 도보로 살펴보아야 한다.
  \item 파크페르메\\
    조직위원회의 결정에 따라서 내구레이싱 경기를 완료한 차량들은 일정 시간 동안 특정 위치에 보관되어 검차를 진행할 수 있으며, 이 시간 동안 임의 이동 및 조작, 정비 등의 행위는 금지된다. 또한 검차 중 규정위반 사항이 적발될 경우 실격시킬 수 있다.
\end{enumerate}

\section{에너지 효율성}
\begin{enumerate}
  \item 내구레이스 주행을 완료하고 파크페르메에 입고된 차량에 한하여 연비/전비를 측정한다.
  \item 에너지 효율성은 측정된 연료 소비량 또는 에너지 사용량과 참가 차량의 내구레이스 평균 랩타임을 기준으로 한다.
  \item 에너지 효율성 경기의 진행(C-Formula)
    \begin{enumerate}
      \item 내구레이스 주행 전 주입연료량 확인용 투명 호스로 식별 가능할 만큼 연료를 채워야 한다.
      \item 주유는 오피셜에 의해 단 한 차례만 이루어지며, 주유 중 차량을 흔들거나 기울이는 등 연료주입에 영향을 끼치는 행동은 금지된다.
      \item 주유를 마치고 연료 주입구 및 연료량 확인 호스에는 봉인 스티커를 붙이며, 이를 고의적으로 훼손하거나 부정행위로 판단될 경우 실격 처리한다.
      \item 봉인 훼손이 필요하거나 의도하지 않은 봉인 훼손이 발생될 경우 오피셜의 판단에 따라 재주유 및 재봉인을 받아야 한다.
      \item 최초 주유 및 봉인 이후 예열, 테스트, 봉인 파손, 연료수위 감소 등으로 인한 재주유 및 재봉인 시 추가되는 연료의 양은 사용 연료량의 2배로 가산한다.
    \end{enumerate}
  \item 에너지 효율성 경기의 진행(E-Formula)
    \begin{enumerate}
      \item 내구레이스 이전 차량의 배터리는 완전 충전을 권장한다.
      \item 차량기술규정 에너지미터 관련 규정 위반 시 해당 규정에 의해 패널티가 적용된다.
      \item 회생 제동 에너지가 사용 에너지를 초과하는 경우 해당 차량은 에너지 효율성 점수에 만점을 부여한다.
    \end{enumerate}
  \item 에너지 효율성 점수계산
    \begin{enumerate}
      \item 보정 에너지 소비량\\
        각 연료 및 에너지는 다음 식에 따라 보정 에너지 소비량으로 변환된다. 단, 휘발유의 계산공식은 조직위원회의 판단에 따라 부피 혹은 질량 단위 중 하나로 계산된다.
        \[
          1\,\mathrm{L\ Gasoline} = 2.31\,\mathrm{kg\ of\ CO_2}
          \quad\text{or}\quad
          1\,\mathrm{kg\ Gasoline} = 2.95\,\mathrm{kg\ of\ CO_2}
        \]
        \[
          1\,\mathrm{kWh\ Electric} = 1.3\,\mathrm{kg\ of\ CO_2}
        \]
      \item E-Formula의 회생제동 충전에 의한 에너지 소비량은 다음 식과 같이 계산한다.
        \[
          E_{\mathrm{calculated}} = E_{\mathrm{used}} - E_{\mathrm{charged}}
        \]
      \item 경기 점수 계산\\
        $CO_{2\min}$은 효율성 이벤트에서 실격하지 않은 차량 중 가장 적은 양의 랩당 평균 보정 에너지 소비량, $CO_{2\mathrm{yours}}$는 해당 팀의 랩당 평균 보정 에너지 소비량, $T_{\min}$은 효율성 이벤트에서 실격하지 않은 차량 중 가장 빠른 평균 랩타임, $T_{\mathrm{yours}}$는 해당 팀의 평균 랩타임이다.
        \[
          EF = \frac{T_{\min}}{T_{\mathrm{yours}}}
               \times \frac{CO_{2\min}}{CO_{2\mathrm{yours}}}
        \]
        \[
          \text{에너지 효율성 경기점수} = 50 \times
          \frac{(EF_{\min}/EF_{\mathrm{yours}})-1}
               {(EF_{\min}/EF_{\max})-1}
        \]
    \end{enumerate}
  \item 에너지 효율성 실격 사항
    \begin{enumerate}
      \item 평균 내구 랩타임이 내구 이벤트를 완주한 가장 빠른 팀의 평균 내구 랩타임의 1.45배를 초과하는 차량은 실격된다.
      \item 보정된 랩당 평균 에너지 소비량이 $60.06\,\mathrm{kg\ of\ CO_2}/100\,\mathrm{km}$를 초과하면 실격된다.
      \item 내구레이스 입차 30분 전까지 주유 및 봉인받지 못한 C-Formula 팀은 실격된다.
      \item 차량 연료탱크 설계 등의 문제로 주유 완료까지 소요되는 시간이 20분 이상인 차량은 실격시킬 수 있다.
      \item 내구레이스에 완주하지 못한 차량은 실격된다.
      \item 위원회의 판단하에 차량의 손상 및 잠재적 위험의 발생이 우려된다면 재주유 및 에너지미터 계측을 실시하지 않으며, 실격시킬 수 있다(예: 연료 누출, 배터리 화재 등).
    \end{enumerate}
\end{enumerate}

\section{채점 및 진행순서}
\begin{enumerate}
  \item Formula 부문의 종목별 채점은 별표 1의 채점기준표에 따른다.
  \item 내구레이스 경기는 오토크로스 경기 성적순으로 진행하는 것을 원칙으로 하되 조직위원회 결정에 따라 변경될 수 있다.
  \item 모든 경기는 차량당 2명의 드라이버가 진행한다.
  \item 각 경기는 대기 장소 입장 마감 시간이 정해질 수 있으며, 마감 시간 전에 대기 장소에 입장하지 않을 경우에는 경기 참가가 제한될 수 있다.
  \item 각 경기는 진행 및 종료 시간이 정해질 수 있으며 종료 시간 전에 경기를 마친 차량에 대해서만 점수가 부여된다.
  \item 최종 합계 점수가 동일할 경우 내구경기, 오토크로스 경기, 보고서 평가, 스키드패드 경기, 가속경기, 에너지 효율성의 순서로 높은 성적을 우선으로 하여 순위를 결정하고, 모든 경기에서의 점수가 동일할 경우 운전자가 탑승하지 않고 연료 및 오일이 가득 찬 상태의 공차 무게가 높은 차량을 우선으로 순위를 결정한다.
\end{enumerate}

\section{날씨 상태}
\begin{enumerate}
  \item 조직위원회는 날씨 상태에 따라 경기별 점수, 진행상황을 조정할 수 있다.
  \item 조직위원회는 날씨 상태에 따라 Dry / Damp / Wet의 3가지 상태를 결정한다.
  \item 날씨 상태와 각 상태별 허가된 타이어는 다음과 같다.
    \begin{enumerate}
      \item Dry: 레인타이어를 제외한 모든 타이어
      \item Damp: 모든 타이어
      \item Wet: 슬릭타이어를 제외한 모든 타이어
    \end{enumerate}
  \item 조직위원회는 날씨 상태에 따라 경기를 연기 또는 중지할 수 있으며, 차량은 즉시 그 지시를 따라야 한다.
\end{enumerate}

\chapter{차량의 정비}

\section{차량 정비}
\begin{enumerate}
  \item 차량의 정비는 패독 내에서만 가능하며, 반드시 구동 바퀴는 지면으로부터 10cm 이상 떨어진 상태에서만 정비할 수 있다.
  \item 차량정비를 위한 패독 내 정비구역은 정비를 위한 인원을 제외한 사람의 출입이 통제될 수 있도록 안전띠 등으로 구분되어야 한다.
  \item 차량정비인원은 정비구역 내에 동시에 최대 5인만 가능하며, 차량정비인원은 지정된 조끼 등으로 구별되어야 한다.
  \item 정비구역 외부에서 공구 등을 전달하거나, 차량의 상하차 등 많은 인원이 필요할 경우는 차량정비인원에 포함되지 않으며, 지정된 차량정비인원에 의해 안전 통제가 되어야 한다.
  \item 차량정비 시 상기 사항을 준수하지 않거나, 안전사고를 유발하는 행위를 비롯하여 안전사고가 발생하는 경우 또는 특히 전기 Formula 고전압 관련 감전 등과 같은 안전사고에 위배된다고 판단될 경우 조직위원회에서 감점을 부가하거나 실격 처리할 수 있다.
\end{enumerate}

\section{주유 및 충전}
\begin{enumerate}
  \item 주유 및 충전은 별도 지정된 구역에서만 가능하다.
  \item 유류의 보관: 별도 지정 공간에 유류를 보관하여야 한다.
  \item 차량의 주유는 연료 주입 속도와 무관하게 연료 탱크나 차량을 어떤 방식으로도 조작하지 않고 용량 끝까지 채울 수 있어야 한다.
  \item 에너지 효율성 경기의 공정성을 위해 오토크로스 주행을 완료한 차량은 내구레이스 진행 전까지 단 1회의 주유 및 봉인이 가능하며, 추가 주유 시 제12조 3항을 적용한다.
\end{enumerate}

\RuleIndexEnd

\chapter*{부칙}
\begin{enumerate}
  \item 이 규정은 제정일로부터 시행한다.
  \item 이 규정의 제정 및 개정 이력은 아래와 같다.
\end{enumerate}

\begin{center}
\begin{tblr}{colspec={Q[l]Q[l]Q[l]Q[l]},rowsep=3pt}
2011.8.1 제정 & 2012.3.8 개정 & 2013.3.8 개정 & 2014.3.13 개정 \\
2016.3.10 개정 & 2017.6.15 개정 & 2018.3.15 개정 & 2019.3.21 개정 \\
2020.3.19 개정 & 2021.3.18 개정 & 2022.4.21 개정 & 2023.4.27 개정 \\
2024.4.11 개정 & 2025.4.18 개정 & 2026.3.10 개정 &
\end{tblr}
\end{center}

\clearpage
\thispagestyle{star}
\begin{center}
  {\pretendardb\Large Formula 부문 종목별 채점기준표}
\end{center}

\begin{center}
\begin{tblr}{
  colspec={Q[l,5cm]Q[c,3cm]},
  hlines,
  vlines,
  row{1}={font=\pretendardb,bg=gray!20},
  cells={valign=m}
}
종목 & 점수 \\
보고서 평가 & 200점 \\
가속 경기 & 100점 \\
스키드패드 경기 & 100점 \\
오토크로스 경기 & 200점 \\
내구레이싱 경기 & 350점 \\
에너지 효율 & 50점 \\
합계 & 1,000점
\end{tblr}
\end{center}

}

\section{안전검사}
\begin{enumerate}
  \item 각 차량이 규정의 요구사항을 따르는지를 결정하기 위하여 검사가 진행된다. 검사는 드라이버의 안전장비 점검과 드라이버 탈출시간 시험을 포함한다. 안전검사를 통과한 후에 틸팅검사, 소음검사/우천검사, 제동검사 순으로 진행되며 대회 운영 상황에 따라 검사 순서는 변경될 수 있다.
  \item 참가팀은 안전검사 시 아래 항목이 준비되어야 한다.
    \begin{itemize}
      \item 차량
      \item 모든 드라이버
      \item 타이어(드라이 \& 레인 타이어)
      \item 공식 인증 및 방염소재의 운전자 보호장비(Formula 차량기술규정 참조: 팔안전벨트, 헬멧, 글러브, 슈트 등)
      \item 소화기
      \item 내연기관동력 차량: 푸쉬바 \& 퀵 잭(Quick Jack)
      \item 전기동력 차량: 절연 푸쉬바 \& 퀵 잭(Quick Jack)
      \item 구조 대응물 양식(Safety Structure Equivalency) \& 사용 재료 물성 증명서 등(사전제출, 일정 별도 공지)
      \item 충돌흡수장치 설계 및 해석 등 설명자료(사전제출, 일정 별도 공지)
      \item 안전검사 요구사항 만족함을 증명하는 문서(필요 시 제출, 조직위에서 협의 후 결과 결정)
      \item 전기동력 차량: Formula 차량기술규정 전기시스템 검사에 해당하는 품목
    \end{itemize}
  \item E-Formula 차량은 Formula 차량기술규정에 의거하여 전기 기술 검사를 진행한다. 전기 기술 검사 중 우천검사는 반드시 안전검사와 틸팅검사를 완료한 후 진행한다.\\
    (E-Formula 차량 검사 순서: 안전검사 $\rightarrow$ 틸팅검사 $\rightarrow$ 우천검사 $\rightarrow$ 음향/제동검사)
\end{enumerate}
\section{틸팅검사}
\begin{enumerate}
  \item Formula 차량기술규정에 의거하여 틸팅 장비로 검사를 진행한다. 검사는 팀의 다이나믹 경기에 참여하는 전체 드라이버 중(예비 드라이버 포함) 키가 가장 큰 드라이버가 보호장비를 모두 착용 후 탑승하여 실행한다.
  \item 틸팅검사는 연료를 가득 주유하고 각종 오일류 및 냉각수를 규정에 따라 채운 상태에서 실행한다.
  \item 만약 연료주입구가 중앙에 위치하지 않을 경우 연료주입구가 있는 방향으로 틸팅 검사를 실행한다.
  \item 기울기 45° 기준 연료 및 액체(오일, 냉각수 등)류 유출을 확인하며, 유출 시 불합격 처리한다.
  \item 기울기 60° 기준 차량 전복 여부를 확인하며, 위쪽 바퀴가 바닥에서 떨어질 시 불합격 처리한다. (횡G 1.7G 상당에서의 전복(roll over) 방지)
\end{enumerate}

\section{소음검사/음향검사/우천검사 및 제동검사}
\begin{enumerate}
  \item 내연기관동력의 차량은 소음측정기를 이용하여 소음검사를 진행한다.
  \item 전기동력 차량은 소음측정기를 이용하여 Formula 차량기술규정에서 규정하는 운전 준비 완료 상태음(RTDS)에 대한 음향 레벨 검사를 진행한다. 이때 측정을 위한 차량 주변 반경 2m의 중심은 RTDS 음원을 기준으로 하고, 측정은 이 중심과 동등한 높이에서 측정된다.
  \item 전기동력의 차량은 Formula 차량기술규정에 의거하여 우천검사를 진행한다. 우천검사를 통과한 차량은 타이어를 교체하거나 타이어를 건조시킨 후 제동검사를 실시한다.
  \item 제동검사는 차량이 주행하는 중에 제동하여 모든 바퀴가 제동되는지를 검사한다.
  \item 제동검사에서는 다음과 같은 항목을 점검한다.
    \begin{itemize}
      \item 모든 바퀴 제동
      \item 제동 직후 차량 회전각도 상태(차체회전각이 45° 이내이어야 할 것)
      \item 제동 완료 후 동력상실 및 자력 재출발 가능 여부
      \item 일정 거리 내 정차 여부
    \end{itemize}
  \item 소음검사/우천검사 및 제동검사 시 검사장 안에 출입 가능 인원은 제한된다.
  \item 회생 제동(Regenerate brake)을 사용하는 E-Formula의 경우, 제동을 시작하기 전에 드라이버는 운전석의 보조 비상 정지 스위치를 작동시켜 차단회로를 개방시켜야 한다. 회생 제동의 영향 없이 유압 제동장치만으로 제동이 가능해야 한다.
\end{enumerate}

\chapter{경기 진행 및 평가}

\section{보고서 평가}
\begin{enumerate}
  \item 모든 참가팀은 대회 전 설계보고서와 비용보고서 및 요청된 보고서를 제출해야 한다.
  \item 보고서 평가는 보고서 제출과 현장평가로 이루어진다.
  \item 보고서 평가 점수는 총 200점으로 마감 일정 전에 제출 시 최저 점수 50점을 부여하며, 제출된 보고서와 현장평가를 통해 150점을 부여한다. 보고서와 현장평가의 세부평가 점수는 조직위원회에서 결정한다.
  \item 현장평가는 대회 기간 중 진행하며 제출된 보고서와 실제 차량을 비교하여 평가한다.
  \item 보고서 양식, 제출 및 현장평가 세부진행방법, 일정 등은 별도 공지한다.
  \item 제출된 보고서 및 발표자료 등은 전체 공개할 수 있다.
\end{enumerate}

\section{가속경기}
\begin{enumerate}
  \item 가속경기는 편평한 직선 포장도로에서 차의 가속도를 평가한다.
  \item 가속경기의 진행\\
    가속주행 구간을 설정하여 출발점에서 정해진 거리를 가속주행 후 체크지점에서 시간을 체크하여 채점기준에 의거 점수를 부여한다.
  \item 가속경기의 점수
    \begin{enumerate}
      \item 다음 공식은 가속 경기 점수를 산출하는데 사용된다.
        \[
          \text{가속 경기 점수} = 75 \times
          \frac{((1.5 \times T_{\min}) / T_{\mathrm{your}}) - 1}
               {((1.5 \times T_{\min}) / T_{\min}) - 1} + 25
        \]
        여기서 $T_{\mathrm{your}}$은 벌칙 시간을 포함한 최고 랩타임, $T_{\min}$은 가장 빠른 차량이 기록한 시간이다.
      \item 가장 빠른 기록의 150\%를 초과한 차량은 가속경기 점수 중 주행성능에 대한 점수는 부여하지 않고 25점의 완주점수만 주어진다. 완주하지 못한 차량에 대해서는 DNF 처리한다.
    \end{enumerate}
  \item 코스 이탈\\
    경기장 좌우 폭은 2.5m로 제한하며, 코스를 이탈한 경우 차량은 경기를 통과하지 못한 것으로 간주되어 DNF(Did Not Finish) 처리한다.
\end{enumerate}

\fig{가속 경기 및 스키드 패드 경기 코스 예시}{competition}{0.75}

\section{스키드패드 경기}
\begin{enumerate}
  \item 스키드패드 경기는 차량이 일정한 반경을 회전할 때의 선회 성능을 평가하는 경기이다.
  \item 스키드패드 경기코스
    \begin{enumerate}
      \item 스키드패드 경기를 위한 기본적인 코스는 8자 모양의 두 개의 원으로 구성된다. 두 원의 중심간의 거리는 18.25m, 내부원의 직경은 15.25m, 외부원의 직경은 21.25m로 이루어진다. 내부 원과 외부 원 사이의 주행구간의 폭은 3m로 이루어진다.
      \item 기록측정이 이루어지는 출발과 도착선은 양쪽 원이 접하는 구간에 있으며, 차량이 출발선상을 통과하고 돌아올 때까지의 시간을 기록한다.
      \item 내부원은 30도 간격으로 12개의 파일런을 놓으며 외부원은 15도 간격으로 18개를 설치한다.
      \item 각 원에는 파일런이 놓여진 위치를 표시한다.
      \item 입구와 출구를 표시하기 위해 추가된 파일런이 놓여지고, 코스 완주를 하지 않고 코스를 벗어나는 것을 막기 위해 출구 쪽 중앙에 파일런이 놓여질 수 있다.
    \end{enumerate}
  \item 스키드패드 경기의 진행
    \begin{enumerate}
      \item 두 개의 원형트랙을 사용할 경우 차량은 그림 2와 같이 입구쪽(IN)으로 진입하여 원하는 원형 트랙 방향으로 두 바퀴, 반대 방향으로 두 바퀴를 선회한 후 입구 반대쪽의 출구쪽(OUT) 종료 라인을 통과하여야 한다.
      \item 최종 주행시간은 두 번째 회전과 네 번째 회전의 계측시간을 합산한다.
      \item 조직위원회에서는 경기 현장 당시의 계측 장비 설치 위치 및 장비 보호 등의 사유로 우선 진입 트랙 방향을 임의로 결정할 수 있다.
    \end{enumerate}
  \item 스키드패드 경기 벌칙
    \begin{enumerate}
      \item 파일런을 넘어뜨렸거나 밖으로 주행했을 때\\
        게이트의 파일런을 포함하여 모든 파일런을 넘어뜨리거나 파일런 밖으로 주행했을 때는 개당 0.3초의 벌칙 시간이 부과된다.
      \item 코스 이탈\\
        차량은 코스를 완전히 이탈하지 않은 상태에서 연속적인 스핀아웃은 가능하다. 하지만 차량의 4바퀴가 코스를 완전히 이탈하면 DNF 처리한다.
      \item 부정확한 주행\\
        주어진 주행 횟수를 모두 완료하지 못한 차량은 DNF 처리한다.
    \end{enumerate}
  \item 스키드패드 경기점수
    \begin{enumerate}
      \item 스키드패드 경기의 기록은 2번째와 4번째 회전의 평균으로 측정(즉, 원형트랙의 왼쪽 1바퀴와 원형 트랙의 오른쪽 1바퀴의 평균)되며, 벌칙 시간을 포함하여 기록된다.
      \item 스키드패드 경기의 점수는 2번째와 4번째 회전의 측정기록 평균과 벌칙 시간을 포함한 주행시간을 사용하여 계산한다.
      \item 벌칙 시간은 기록이 측정되는 회전 수에 상관없이 전 회전수에서 발생하는 기록을 집계하여 포함한다.
      \item 다음 공식은 스키드패드 점수를 산출하는데 사용된다.
        \[
          \text{스키드패드 점수} = 75 \times
          \frac{((1.25 \times T_{\min}) / T_{\mathrm{your}}) - 1}
               {((1.25 \times T_{\min}) / T_{\min}) - 1} + 25
        \]
        $T_{\mathrm{your}}$은 벌칙시간을 포함한 최고 주행시간(주행 시 왼쪽, 오른쪽 주행시간의 평균값), $T_{\min}$은 참가차량 중 가장 빠른 주행시간이다.
    \end{enumerate}
  \item 가장 빠른 기록의 125\%를 초과한 차량은 스키드패드 경기 점수 중 주행성능에 대한 점수는 부여하지 않고 25점의 완주점수만 주어진다. 완주하지 못한 차량에 대해서는 DNF 처리한다.
\end{enumerate}

\fig{스키드 패드 경기 코스}{competition}{0.62}

\section{오토크로스 경기}
\begin{enumerate}
  \item 오토크로스 경기는 주어진 코스에서 차량의 가속도, 제동 그리고 코너링 성능의 특성을 평가하기 위해 조합된 코스로 만들어진다.
  \item 오토크로스 경기코스
    \begin{enumerate}
      \item 경기코스는 직선구간, 회전구간, 슬라럼구간, 시케인구간 등을 포함하여 구성된다.
      \item 1랩의 길이는 약 700--900m이며, 다음 기준에 따라 구성된다.
        \begin{itemize}
          \item 직선 구간: 양쪽 끝에 헤어핀을 가진 최대 70m 구간 또는 끝부분에 완만한 곡선을 가진 60m를 넘지 않는 구간. 직선구간에는 추월구간이 포함되어 있을 수 있다.
          \item 회전 구간: 30m에서 60m의 직경을 가진 구간
          \item 헤어핀 구간: 최소 9m의 외부 직경을 가진 구간
          \item 슬라럼(지그재그) 구간: 7m에서 15m 사이의 직선으로 놓여진 파일런 구간
          \item 기타 구간: 시케인(자동차 경주도로에서 감속을 위한 장애물), 다양한 회전, 반경이 줄어드는 회전 구간 등을 가지며 트랙의 최소 폭은 3.5m로 만들어진다.
        \end{itemize}
      \item 오토크로스 경기코스는 내구레이싱 경기코스와 유사한 형태로 만들어지며, 주행방향은 역방향으로 진행한다.
      \item 경기장 조건과 상황, 안전 사항 등을 고려하여 경기코스 형태는 다양한 형태로 변형될 수 있다.
    \end{enumerate}
  \item 오토크로스 경기의 진행
    \begin{enumerate}
      \item 차량이 계측선상을 통과하는 순간 또는 신호등에 녹색등이 점등되는 순간 계측을 시작하고 주행을 완료하여 정차구간의 계측선상을 통과할 때까지의 기록을 측정하며, 계측 시작 방식은 조직위원회에서 경기장 상황에 따라 결정한다.
      \item 경기는 두 개의 그룹으로 나누어 진행하며, 각 그룹에는 한 명의 드라이버만 참여가 가능하다.
      \item 드라이버 1인당 2번의 기회가 주어지며, 총 4회의 주행이 가능하다.
      \item 각 그룹에서 주행 우선권은 경기장 입장순으로 하되, 그룹 경기를 처음 참여하는 드라이버에게 우선권을 부여한다.
      \item 조직위원회는 날씨 또는 기술적인 지연 등으로 인한 경기 진행에 대해 조정할 수 있으며, 경기 진행에 대한 제한 시간을 결정할 수 있다. 만약 제한된 시간 안에 주행을 마치지 못한 차량은 DNF 처리한다.
    \end{enumerate}
  \item 오토크로스 경기벌칙\\
    주행이 완료된 차량은 주행시간에 벌칙별로 아래의 시간이 추가된다.
    \begin{enumerate}
      \item 파일런이 넘어지거나 움직인 경우\\
        파일런 하나당 2초가 최종시간에 부과된다.
      \item 코스 이탈\\
        코스 이탈은 차량의 네 바퀴 모두가 파일런이나 코스 선 등 정해진 코스의 경계선을 넘어설 경우를 말한다. 코스를 이탈하게 되면 드라이버는 이탈 발생지점이나 발생지점 전에서 코스로 재진입하여야 하며, 이를 위반 시 20초의 벌칙 시간이 더해진다. 코스 이탈 중 지나친 파일런 하나당 2초의 벌칙 시간이 추가로 부과된다.
      \item 슬라럼을 놓친 경우\\
        주어진 슬라럼 파일런의 하나 또는 그 이상의 게이트를 놓친 경우 각각 하나의 코스 이탈로 계산되어 각 20초의 벌칙시간이 부과된다.
      \item 주행정지 또는 운행불가 차량\\
        차량이 주행 중 멈추거나 외부의 도움 없이 재출발이 불가능할 경우 운행 불가 차량으로 간주되며 DNF 처리한다.
    \end{enumerate}
  \item 오토크로스 주행기록시간\\
    주행기록시간은 주행 랩타임과 주어진 벌칙시간을 합쳐 기록된다.
  \item 오토크로스 경기 점수\\
    오토크로스 점수를 산정하기 위해 다음의 공식이 사용된다.
    \[
      \text{오토크로스 점수} = 150 \times
      \frac{((1.45 \times T_{\min}) / T_{\mathrm{your}}) - 1}
           {((1.45 \times T_{\min}) / T_{\min}) - 1} + 50
    \]
    $T_{\mathrm{your}}$은 벌칙 시간을 포함한 최고 주행시간, $T_{\min}$은 참가 차량 중 가장 빠른 주행시간이다.
  \item 가장 빠른 기록의 145\%를 초과한 차량은 오토크로스 경기 점수 중 주행성능에 대한 점수는 부여하지 않고 50점의 완주 점수만 주어진다. 완주하지 못한 차량에 대해서는 DNF 처리한다.
\end{enumerate}

\CompetitionRulesTail

\end{document}
